\documentclass[12pt]{article}
\usepackage[utf8]{inputenc}
\usepackage[T1]{fontenc}
\usepackage[margin=1in]{geometry}
\usepackage{setspace}
\usepackage{xcolor}
\usepackage[colorlinks=true, urlcolor=blue!60!black, linkcolor=blue!60!black]{hyperref}
\setlength{\parindent}{2em}
\setlength{\parskip}{1.25em}
\renewcommand{\baselinestretch}{1.0}
\begin{document}
\begin{center}
{\LARGE \textit{Frontiero v. Richardson}, 411 U.S. 677 (1973)\par}
\vspace{0.8em}
{\large Case Brief\par}
\end{center}
\vspace{1em}

\section{Memory Jogger}
In an 8--1 judgment, the Supreme Court invalidated federal military benefit statutes that automatically treated servicemen's wives as dependents while requiring servicewomen to prove their husbands' actual dependency, with a four-Justice plurality famously declaring sex a suspect classification subject to strict scrutiny (\href{https://www.courtlistener.com/opinion/108781/frontiero-v-richardson/}{CourtListener}; \href{https://supreme.justia.com/cases/federal/us/411/677/}{Justia}).

\section{Detailed Case Facts}
Sharron Frontiero, a lieutenant in the United States Air Force stationed at Maxwell Air Force Base in Montgomery, Alabama, applied for increased ``basic allowance for quarters'' (housing benefits) and comprehensive medical and dental benefits for her husband, Joseph Frontiero, as a ``dependent'' (\href{https://supreme.justia.com/cases/federal/us/411/677/}{Justia}; \href{https://www.oyez.org/cases/1972/71-1694}{Oyez}). Under the applicable federal statutes---37 U.S.C.~\S\S~401 and 403 (originally enacted in the Career Compensation Act of 1949) and 10 U.S.C.~\S\S~1072 and 1076 (enacted in the Dependents' Medical Care Act of 1956)---Congress established a fringe-benefit scheme to attract career military personnel through reenlistment (\href{https://supreme.justia.com/cases/federal/us/411/677/}{Justia}).

The statutes drew an explicit sex-based line in defining ``dependent.'' Under 37 U.S.C.~\S~401(1) and 10 U.S.C.~\S~1072(2)(A), a male member of the uniformed services could claim his wife as a dependent automatically, without showing that she relied on him for any portion of her financial support (\href{https://supreme.justia.com/cases/federal/us/411/677/}{Justia}). By contrast, under 37 U.S.C.~\S~401 and 10 U.S.C.~\S~1072(2)(C), a female member of the uniformed services could not claim her husband as a dependent unless she proved that he was in fact dependent on her for more than one-half of his support (\href{https://supreme.justia.com/cases/federal/us/411/677/}{Justia}).

Joseph Frontiero was a full-time college student at Huntingdon College in Montgomery, Alabama. According to the parties' stipulated facts, his monthly living expenses (including his share of household expenses) totaled approximately \$354, and he received \$205 per month in veterans' benefits (\href{https://supreme.justia.com/cases/federal/us/411/677/}{Justia}). Because his veterans' benefits covered more than half of his living expenses, Lieutenant Frontiero could not show that he depended on her for over one-half of his support, and the Air Force denied her application. Had she been a male officer married to a wife with identical income and expenses, the benefits would have been granted automatically (\href{https://supreme.justia.com/cases/federal/us/411/677/}{Justia}).

\section{Procedural History}
Sharron and Joseph Frontiero sued Secretary of Defense Melvin Laird (later succeeded on appeal by Secretary of Defense Elliot Richardson) in the United States District Court for the Middle District of Alabama, seeking a permanent injunction against enforcement of the discriminatory statutory provisions and an order granting Lieutenant Frontiero the same housing and medical benefits afforded to similarly situated servicemen (\href{https://supreme.justia.com/cases/federal/us/411/677/}{Justia}; \href{https://law.justia.com/cases/federal/district-courts/FSupp/341/201/1456940/}{Justia District Court}). Because the suit challenged the constitutionality of federal statutes, a three-judge district court was convened, consisting of Circuit Judge Richard T. Rives and District Judges Frank M. Johnson, Jr., and Frank H. McFadden (\href{https://law.justia.com/cases/federal/district-courts/FSupp/341/201/1456940/}{Justia District Court}).

On April 5, 1972, the three-judge District Court, in a 2--1 decision written by Judge Rives, upheld the statutes under traditional rational-basis review (\textit{Frontiero v. Laird}, 341 F. Supp. 201 (M.D. Ala. 1972)) (\href{https://law.justia.com/cases/federal/district-courts/FSupp/341/201/1456940/}{Justia District Court}). The District Court majority surmised that because husbands were typically family breadwinners and roughly 99\% of uniformed service members were male, Congress could reasonably presume wives' dependency while requiring proof for husbands to achieve a ``considerable saving of administrative expense and manpower'' (\href{https://supreme.justia.com/cases/federal/us/411/677/}{Justia}). Judge Johnson dissented, arguing that under \textit{Reed v. Reed}, 404 U.S. 71 (1971), administrative convenience could not justify denying benefits to similarly situated women (\href{https://law.justia.com/cases/federal/district-courts/FSupp/341/201/1456940/}{Justia District Court}).

The Frontieros took a direct appeal to the U.S. Supreme Court, which noted probable jurisdiction on October 10, 1972 (409 U.S. 840) (\href{https://supreme.justia.com/cases/federal/us/411/677/}{Justia}). Oral argument was held on January 17, 1973, with Joseph J. Levin, Jr., arguing for the appellants, Samuel Huntington arguing for the government, and Ruth Bader Ginsburg making her first Supreme Court oral argument as \textit{amicus curiae} on behalf of the American Civil Liberties Union (\href{https://www.oyez.org/cases/1972/71-1694}{Oyez}; \href{https://www.law.georgetown.edu/georgetown-law-journal/wp-content/uploads/sites/26/2022/09/Suk-Justice-Ginsburgs-Cautious-Legacy.pdf}{Georgetown Law Journal}). On May 14, 1973, the Supreme Court reversed in an 8--1 judgment (\href{https://supreme.justia.com/cases/federal/us/411/677/}{Justia}).

\section{Judicial Votes}
\begin{itemize}
    \item \textbf{Overall Vote:} 8--1 to reverse the District Court (no majority opinion on the standard of review) (\href{https://supreme.justia.com/cases/federal/us/411/677/}{Justia}; \href{https://www.oyez.org/cases/1972/71-1694}{Oyez}).
    \item \textbf{Plurality Opinion (4 Justices):} Justice William J. Brennan, Jr., announced the judgment of the Court and delivered an opinion joined by Justices William O. Douglas, Byron R. White, and Thurgood Marshall (\href{https://supreme.justia.com/cases/federal/us/411/677/}{Justia}).
    \item \textbf{Concurring in the Judgment (1 Justice):} Justice Potter Stewart filed a brief statement concurring in the judgment (\href{https://supreme.justia.com/cases/federal/us/411/677/}{Justia}).
    \item \textbf{Concurring in the Judgment (3 Justices):} Justice Lewis F. Powell, Jr., filed an opinion concurring in the judgment, joined by Chief Justice Warren E. Burger and Justice Harry A. Blackmun (\href{https://supreme.justia.com/cases/federal/us/411/677/}{Justia}).
    \item \textbf{Dissenting (1 Justice):} Justice William H. Rehnquist filed a one-sentence dissenting statement adopting Judge Rives's District Court opinion (\href{https://supreme.justia.com/cases/federal/us/411/677/}{Justia}).
\end{itemize}

\section{Holding}
\begin{itemize}
    \item \textbf{Judgment:} Reversed. The challenged provisions of 37 U.S.C.~\S\S~401, 403 and 10 U.S.C.~\S\S~1072, 1076 violate the Due Process Clause of the Fifth Amendment insofar as they require a female member of the uniformed services to prove the actual dependency of her spouse to obtain housing and medical benefits that male members receive automatically (\href{https://supreme.justia.com/cases/federal/us/411/677/}{Justia}).
    \item \textbf{Rule Announced:} Because no five Justices agreed on a single standard of review, the binding rule of \textit{Frontiero} is the narrowest ground supporting the judgment: under the equal protection component of the Fifth Amendment's Due Process Clause (via \textit{Bolling v. Sharpe}, 347 U.S. 497 (1954)) and \textit{Reed v. Reed}, 404 U.S. 71 (1971), federal statutes that accord dissimilar treatment to similarly situated male and female service members solely for ``administrative convenience'' are unconstitutional (\href{https://supreme.justia.com/cases/federal/us/411/677/}{Justia}).
    \item \textbf{Holding vs. Dicta:} The narrow holding is that the military's asymmetric spousal-dependency scheme is invalid and that female service members are entitled to the same automatic spousal benefits as male service members (rather than striking down the benefit program entirely) (\href{https://supreme.justia.com/cases/federal/us/411/677/}{Justia}). Justice Brennan's declaration that sex-based classifications are ``inherently suspect'' and trigger ``strict judicial scrutiny'' commanded only four votes and therefore was \emph{not} a holding of the Court; three years later, in \textit{Craig v. Boren}, 429 U.S. 190 (1976), the Court settled on intermediate scrutiny for sex classifications rather than strict scrutiny (\href{https://supreme.justia.com/cases/federal/us/411/677/}{Justia}; \href{https://casssunstein.substack.com/p/the-sex-discrimination-cases-of-the}{Sunstein}). Likewise, Justice Brennan's observations about intelligence and physical disability as ``nonsuspect statuses'' and his commentary on the pending Equal Rights Amendment were non-binding plurality dicta (\href{https://supreme.justia.com/cases/federal/us/411/677/}{Justia}).
\end{itemize}

\section{Analysis of Opinions}

\subsection{Plurality Opinion (Justice Brennan, joined by Justices Douglas, White, and Marshall)}
Justice Brennan began by noting that although the Fifth Amendment lacks an express Equal Protection Clause, its Due Process Clause forbids federal discrimination ``so unjustifiable as to be violative of due process'' (\href{https://supreme.justia.com/cases/federal/us/411/677/}{Justia}). He highlighted the twofold discriminatory impact of the statutes: procedurally, servicewomen had to prove their spouses' dependency while servicemen did not; substantively, servicemen who supplied less than half of their wives' support received benefits while similarly situated servicewomen were denied them (\href{https://supreme.justia.com/cases/federal/us/411/677/}{Justia}).

Reading \textit{Reed v. Reed} as an implicit departure from traditional rational-basis review, Justice Brennan articulated a multi-factor framework to justify treating sex as a suspect classification comparable to race, alienage, and national origin: (1) America's ``long and unfortunate history of sex discrimination,'' rooted in ``romantic paternalism'' exemplified by Justice Bradley's concurrence in \textit{Bradwell v. State}, 83 U.S. (16 Wall.) 130, 141 (1873), which denied women voting, property, and occupational rights; (2) the high visibility of sex and ongoing underrepresentation of women in political institutions (noting in footnote 17 that no woman had ever served as President or Supreme Court Justice, none then sat in the Senate, and only 14 served in the House); (3) the immutability of sex as an ``accident of birth'' unrelated to individual responsibility; (4) the fact that sex ``frequently bears no relation to ability to perform or contribute to society,'' distinguishing it from intelligence or physical disability; and (5) recent congressional consensus against sex discrimination in the Equal Pay Act of 1963, Title VII of the Civil Rights Act of 1964, and Congress's 1972 passage of the Equal Rights Amendment (\href{https://supreme.justia.com/cases/federal/us/411/677/}{Justia}).

Applying strict scrutiny, Justice Brennan rejected the government's sole justification of ``administrative convenience.'' First, the government offered no concrete evidence that granting automatic benefits to all servicemen's wives was actually cheaper than verifying dependency via simple affidavits, especially given census and labor statistics showing that 41.5\% of married women were employed and many military wives would fail an actual dependency test (\href{https://supreme.justia.com/cases/federal/us/411/677/}{Justia}). Second, even if administrative savings existed, ``the Constitution recognizes higher values than speed and efficiency,'' and administrative convenience cannot justify arbitrary sex-based lines (\href{https://supreme.justia.com/cases/federal/us/411/677/}{Justia}).

\subsection{Concurrence in the Judgment (Justice Stewart)}
Justice Stewart issued a one-sentence concurrence agreeing that the challenged statutes ``work an invidious discrimination in violation of the Constitution'' under \textit{Reed v. Reed}, 404 U.S. 71 (1971), without endorsing strict scrutiny or suspect-class status for sex (\href{https://supreme.justia.com/cases/federal/us/411/677/}{Justia}).

\subsection{Concurrence in the Judgment (Justice Powell, joined by Chief Justice Burger and Justice Blackmun)}
Justice Powell agreed that the statutes violated the Fifth Amendment's Due Process Clause, but declined to join Justice Brennan's declaration that sex is an inherently suspect classification (\href{https://supreme.justia.com/cases/federal/us/411/677/}{Justia}). He advanced two reasons for judicial restraint. First, \textit{Reed v. Reed} ``abundantly supports our decision today'' without needing to expand the narrow category of suspect classifications (\href{https://supreme.justia.com/cases/federal/us/411/677/}{Justia}). Second, and most compellingly in his view, Congress had just submitted the Equal Rights Amendment (ERA) to the states for ratification. If adopted, the ERA would resolve the constitutional standard for sex classifications through the Article V democratic process; by declaring sex a suspect class while state legislatures were actively debating ratification, the Court would prematurely preempt a major political decision and weaken confidence in judicial restraint (\href{https://supreme.justia.com/cases/federal/us/411/677/}{Justia}).

\subsection{Dissenting Statement (Justice Rehnquist)}
Justice Rehnquist dissented in a single sentence, stating that he dissented ``for the reasons stated by Judge Rives in his opinion for the District Court, \textit{Frontiero v. Laird}, 341 F. Supp. 201 (1972)'' (\href{https://supreme.justia.com/cases/federal/us/411/677/}{Justia}). By incorporating Judge Rives's opinion, Justice Rehnquist endorsed deferential rational-basis review and concluded that Congress could rationally rely on empirical generalizations about male breadwinners and the overwhelmingly male composition of the military to save administrative costs (\href{https://law.justia.com/cases/federal/district-courts/FSupp/341/201/1456940/}{Justia District Court}).

\section{Examples: Future Applications}
\begin{enumerate}
    \item \textbf{Same-Side (Unconstitutional): Asymmetric Survivors' Benefits in Social Security.} A federal statute automatically pays Social Security survivor benefits to widows of deceased male workers, but requires widowers of deceased female workers to prove they received at least half their support from their wives. Like \textit{Frontiero}, this law rests on archaic breadwinner stereotypes and mere administrative convenience, imposing both a procedural proof burden and substantive benefit denial on female wage earners' families; it fails equal protection review (mirroring \textit{Califano v. Goldfarb}, 430 U.S. 199 (1977)).
    \item \textbf{Same-Side (Unconstitutional): Automatic Parental Leave Presumption for Mothers Only.} A federal agency grants six weeks of paid caregiver leave automatically to female employees upon the birth of a child, while requiring male employees to submit affidavits proving they are the ``primary caregiver'' to receive the same leave. Because the rule treats similarly situated male and female employees differently solely on the basis of sex stereotypes and administrative convenience, it is invalid under \textit{Frontiero} and \textit{Reed}.
    \item \textbf{Opposite-Side (Constitutional): Uniform Income-Based Dependency Test.} Congress amends 37 U.S.C.~\S~401 so that \emph{all} service members---male and female alike---must submit a brief annual tax-return affidavit showing that their spouse earns less than half of the household's living expenses in order to receive dependent housing and medical allowances. Because the classification turns on actual financial need rather than sex, it is facially neutral and easily survives rational-basis review.
    \item \textbf{Opposite-Side (Constitutional): Sex-Neutral Disability Accommodation Standard.} A state statute requires applicants for specialized vocational-rehabilitation subsidies to submit medical documentation of a qualifying physical disability, while waiving documentation for applicants already certified as disabled by the Social Security Administration. As Justice Brennan's plurality noted, physical disability---unlike sex---is directly relevant to an individual's functional needs in this context, and relying on an existing neutral administrative certification does not employ a suspect or sex-based classification (\href{https://supreme.justia.com/cases/federal/us/411/677/}{Justia}).
    \item \textbf{Fence-Sitter (Hard Case): Targeted Remedial Tax Exemption for Widows.} A state grants a modest annual property-tax exemption to widows, but not to widowers, citing statistical evidence that older women face steeper historical wage disparities and greater financial vulnerability after a spouse's death. Under \textit{Frontiero}'s footnote 22, Justice Brennan explicitly distinguished statutes designed ``to rectify the effects of past discrimination against women'' from laws that heap additional economic disadvantages onto women (\href{https://supreme.justia.com/cases/federal/us/411/677/}{Justia}). Under Brennan's strict scrutiny, such a law would still face heavy skepticism unless narrowly tailored, whereas under a more deferential remedial approach (as in \textit{Kahn v. Shevin}, 416 U.S. 351 (1974)), the Court might uphold it as compensating for past economic inequality.
\end{enumerate}

\section{Critique}
\begin{itemize}
    \item \textbf{Originalist and Textualist Critiques:} Originalist scholars and jurists critique \textit{Frontiero} on two distinct textual and historical grounds. First, \textit{Frontiero} relies on ``reverse incorporation'' through the Fifth Amendment's Due Process Clause (\textit{Bolling v. Sharpe}, 347 U.S. 497 (1954)), even though the Fifth Amendment, ratified in 1791, contains no Equal Protection Clause and was not originally understood to forbid sex-based federal benefit rules (\href{https://casssunstein.substack.com/p/the-sex-discrimination-cases-of-the}{Sunstein}; \href{https://academiccommons.columbia.edu/doi/10.7916/3jrg-8056/download}{Columbia Academic Commons}). Second, originalists note that neither the Fifth Amendment nor the Fourteenth Amendment (ratified in 1868) was originally understood to make sex a suspect classification; Justice Brennan's plurality explicitly treated American history and tradition as a cautionary record of ``romantic paternalism'' to be overcome rather than a binding original meaning (\href{https://casssunstein.substack.com/p/the-sex-discrimination-cases-of-the}{Sunstein}; \href{https://constitutioncenter.org/the-constitution/white-papers/democratic-constitutionalism}{National Constitution Center}). From this perspective, Justice Powell's deference to the Article V amendment process was institutionally sound, because elevating sex to a suspect class belongs to the people via constitutional amendment rather than judicial decree.
    \item \textbf{Progressive, Feminist, and Doctrinal Critiques:} Progressive and feminist scholars---including commentators examining Ruth Bader Ginsburg's litigation strategy---praise Justice Brennan's sharp dismantling of ``romantic paternalism'' and his recognition that seemingly benign stereotypes trap women in subordinate roles (\href{https://columbialawreview.org/content/goodbye-old-friend-tribute-to-justice-ruth-bader-ginsburg/}{Columbia Law Review}; \href{https://www.law.georgetown.edu/georgetown-law-journal/wp-content/uploads/sites/26/2022/09/Suk-Justice-Ginsburgs-Cautious-Legacy.pdf}{Georgetown Law Journal}). At the same time, feminist theorists criticize \textit{Frontiero}'s fractured outcome: by falling one vote short of strict scrutiny due to Justice Powell's wait-and-see posture on the ERA---which ultimately failed to achieve ratification within its deadline---the Court left sex discrimination with a weaker tier of ``intermediate scrutiny'' in \textit{Craig v. Boren} (1976) (\href{https://law.yale.edu/sites/default/files/documents/pdf/Faculty/Siegel_SheThePeople.pdf}{Siegel, Yale Law School}; \href{https://www.culawreview.org/journal/bringing-sex-discrimination-under-strict-scrutiny-the-need-for-an-equal-rights-amendment}{Columbia Undergraduate Law Review}). Scholars such as Reva Siegel also observe that the resulting intermediate-scrutiny compromise obscured the deeper constitutional foundations of sex equality, including the Nineteenth Amendment's structural transformation of women's citizenship (\href{https://law.yale.edu/sites/default/files/documents/pdf/Faculty/Siegel_SheThePeople.pdf}{Siegel, Yale Law School}). Additionally, disability-rights scholars critique Justice Brennan's casual dicta dismissing ``physical disability'' and ``intelligence'' as ``nonsuspect statuses'' (\href{https://supreme.justia.com/cases/federal/us/411/677/}{Justia}).
    \item \textbf{Underlying Competing Values:} \textit{Frontiero} exposes a classic tension between \emph{substantive individual equality} (insisting that the government judge people as individuals rather than by immutable group stereotypes) and \emph{democratic and administrative restraint} (deferring to Congress's use of low-cost empirical proxies in benefit programs and respecting the Article V amendment process when fundamental constitutional change is underway).
\end{itemize}

\section{Key Quotations}

\begin{quote}
``There can be no doubt that our Nation has had a long and unfortunate history of sex discrimination. Traditionally, such discrimination was rationalized by an attitude of `romantic paternalism' which, in practical effect, put women not on a pedestal, but in a cage.''
\end{quote}
\noindent --- Justice William J. Brennan, Jr., Plurality Opinion, 411 U.S. at 684 (\href{https://supreme.justia.com/cases/federal/us/411/677/}{Justia}; \href{https://www.courtlistener.com/opinion/108781/frontiero-v-richardson/}{CourtListener}).

\begin{quote}
``And what differentiates sex from such nonsuspect statuses as intelligence or physical disability, and aligns it with the recognized suspect criteria, is that the sex characteristic frequently bears no relation to ability to perform or contribute to society.''
\end{quote}
\noindent --- Justice William J. Brennan, Jr., Plurality Opinion, 411 U.S. at 686--687 (\href{https://supreme.justia.com/cases/federal/us/411/677/}{Justia}).

\begin{quote}
``With these considerations in mind, we can only conclude that classifications based upon sex, like classifications based upon race, alienage, or national origin, are inherently suspect, and must therefore be subjected to strict judicial scrutiny.''
\end{quote}
\noindent --- Justice William J. Brennan, Jr., Plurality Opinion, 411 U.S. at 688 (\href{https://supreme.justia.com/cases/federal/us/411/677/}{Justia}).

\begin{quote}
``And when we enter the realm of `strict judicial scrutiny,' there can be no doubt that `administrative convenience' is not a shibboleth, the mere recitation of which dictates constitutionality.''
\end{quote}
\noindent --- Justice William J. Brennan, Jr., Plurality Opinion, 411 U.S. at 690 (\href{https://supreme.justia.com/cases/federal/us/411/677/}{Justia}).

\begin{quote}
``By acting prematurely and unnecessarily, as I view it, the Court has assumed a decisional responsibility at the very time when state legislatures, functioning within the traditional democratic process, are debating the proposed Amendment.''
\end{quote}
\noindent --- Justice Lewis F. Powell, Jr., Concurring in the Judgment, 411 U.S. at 692 (\href{https://supreme.justia.com/cases/federal/us/411/677/}{Justia}).

\end{document}
